\chapter*{Introduction générale}
\addcontentsline{toc}{chapter}{Introduction générale}

La transformation numérique des restaurants ne se limite plus à la présence en ligne ou à la prise de réservation. Elle touche désormais l'organisation opérationnelle : coordination entre la salle et la cuisine, disponibilité des plats, suivi des stocks, gestion des paiements, fidélisation et analyse des retours clients. Dans un établissement où plusieurs acteurs travaillent en parallèle, une information mal transmise peut provoquer un retard, une erreur de préparation, une rupture non anticipée ou une expérience client dégradée.

Le projet Tastify s'inscrit dans ce contexte. Il vise à proposer une application web centralisée pour la gestion d'un restaurant, en couvrant à la fois les besoins internes du personnel et les interactions côté client. Le dépôt montre une architecture full-stack composée d'un backend Django REST Framework, de services Redis et Celery, d'une base MySQL, de WebSocket via Django Channels et de deux applications React : une application staff et un portail client.

La particularité du projet réside dans son périmètre transversal. Tastify ne se limite pas à un module de menu ou à une simple caisse : il relie les tables, les commandes, le KDS, les stocks, les réservations, les paiements, la fidélité, les avis clients et les tableaux de bord. Cette intégration rend le projet pertinent pour un rapport de Projet de Fin d'Année, car elle mobilise plusieurs dimensions du génie logiciel : modélisation, architecture, API, sécurité, temps réel, asynchronisme, expérience utilisateur et validation.

Le rapport suit la structure de l'ancien rapport fourni, sans réutiliser son domaine métier. Le premier chapitre présente le contexte général, la problématique, l'étude de l'existant et les objectifs. Le deuxième chapitre formalise l'analyse des besoins et la conception. Un chapitre spécifique est consacré à l'analyse de sentiment, car cette fonctionnalité combine un enjeu métier et un pipeline technique précis. Les chapitres suivants décrivent l'architecture, la réalisation, les interfaces, les tests, les limites et les perspectives.

Un principe de rigueur guide la rédaction : les fonctionnalités sont décrites uniquement lorsqu'elles sont confirmées par le dépôt ou par les fichiers fournis. Les informations incertaines sont signalées explicitement, afin d'éviter de présenter au jury des éléments non vérifiables.
